
\subsection{A division of labor between CPT and SFT}
The recurring shape of our results --- base $<$ BioCPT $>$ SFT on every held-out axis --- points
to a division of labor that is usually conflated. Continued pretraining does not simply add
domain facts; it \emph{re-organizes and lifts a shared capability substrate} that is then
visible on tasks (general knowledge, code) far from the pretraining domain. Supervised
fine-tuning does the opposite: it \emph{narrows} the model onto the target task distribution,
trading breadth on held-out axes for focus on the served task. Neither stage is a substitute
for the other, and budgeting them as interchangeable --- or omitting CPT to ``save compute'' ---
misreads their roles. This is consistent with, and sharpens, prior router-level evidence on this
model family that CPT reshapes middle-layer computation while SFT re-aligns the output.

\subsection{Biological sequence as structured scientific data}
If narrow, out-of-distribution biological sequence \emph{raises} general and coding competence,
then the ``narrow domain = capacity competition'' framing is inadequate. We propose instead that
biological sequence functions as \emph{structured scientific data}: a corpus whose value to a
foundation model lies not (only) in the facts it carries but in the structural regularities it
forces the model to internalize --- regularities that transfer to other structured/formal
languages such as code. Under this view biology is simply the first representative member of a
broader class (code, mathematics, chemistry/SMILES), and the natural next question is not
``does biology help biology'' but ``how does the \emph{composition} of scientific-data CPT
shape the capability profile of a foundation model.''

\subsection{Mixture ratio is a controllable, general-preserving lever}
The controlled sweep of \S\ref{sec:phase2} answers that question for this model. The two axes do
not trade off: across a biological share of $0\!-\!50\%$ the general axis is flat while the
biology axis rises monotonically, so there is no ``forgetting threshold'' to budget around
within this range. Equally, the composition ablation shows the lever has structure --- DNA and
protein amplify \emph{different} biological capabilities --- so a practitioner can, in principle,
target a desired capability profile by choosing not just how much scientific data to add but
which kind. This is the operational form of the data-centric thesis: training-data composition
is a design surface for the capability profile, not merely a source of facts.

\subsection{Limitations}
This is a training-free re-analysis of one lineage of one 26B MoE model; the CPT / SFT recipe is
fixed, so we characterize \emph{this} lineage rather than establishing a scaling law. Part of the
ARC/HellaSwag gain is a format effect (instruct-vs-completion loglikelihood), which we flag and
which does not affect the knowledge-heavy MMLU gain. Zero-shot sequence-homology is not
measurable by simple probes, so that axis is argued indirectly. TruthfulQA is a genuine, if
small, regression and we do not hide it. The controlled experiment of \S\ref{sec:phase2}
addresses the observational limitation directly --- it fixes the recipe and varies only the data
mixture --- but is itself bounded: each model sees $100$M CPT tokens (a pilot budget chosen to
find trends, not a scaling law), on one 26B MoE model. Whether the general-preserving,
biology-lifting mixture behavior holds at $10\times$ the token budget, across model sizes, and
for other scientific-data domains (code, mathematics, chemistry) remains open, and is the natural
next step now that the lever itself is established.
